% ============================================================================
% CHAPTER 9: RETRIEVAL-AUGMENTED GENERATION
% The dominant paradigm for building knowledge-grounded LLM applications.
% One of the most commonly tested system design topics at FAANG interviews.
% ============================================================================

\chapter{Retrieval-Augmented Generation}
\label{chap:rag}

\begin{tldr}
RAG is the dominant paradigm for building knowledge-grounded LLM applications.
This chapter covers the complete RAG architecture from query processing through
retrieval (sparse, dense, hybrid) to generation, including chunking strategies,
vector databases, re-ranking, and advanced patterns like HyDE and agentic RAG.
RAG system design is one of the most common interview topics for applied ML
roles---every FAANG company building LLM-powered products uses some form of
retrieval augmentation, and interviewers expect candidates to reason about the
full pipeline, not just the model.
\end{tldr}

\bigskip

% ============================================================================
% SECTION 1: RAG ARCHITECTURE OVERVIEW
% ============================================================================
\section{Why RAG and the Core Architecture}
\label{sec:rag-overview}

\subsection{The Problem RAG Solves}

Large language models are remarkably capable, but they suffer from three
fundamental limitations that make them insufficient as standalone knowledge
systems:

\begin{enumerate}
    \item \textbf{Knowledge cutoff.}  An LLM's parametric knowledge is frozen
          at pretraining time.  It cannot answer questions about events,
          documents, or data that appeared after its training data was collected.
    \item \textbf{Hallucination.}  When an LLM encounters a question whose
          answer is not well-represented in its parameters, it generates
          plausible-sounding but factually incorrect text.  There is no
          mechanism to distinguish ``I know this'' from ``I am guessing.''
    \item \textbf{No access to private data.}  Enterprise documents, internal
          wikis, customer records, and proprietary databases are not in the
          pretraining corpus.  Fine-tuning on private data is expensive,
          raises data leakage concerns, and must be repeated as data changes.
\end{enumerate}

\textbf{Retrieval-Augmented Generation} (Lewis et al., NeurIPS 2020) addresses
all three by giving the LLM access to an external knowledge base at inference
time.  Instead of relying solely on parametric memory, the model retrieves
relevant documents and uses them as grounding context for generation.

\subsection{The Core RAG Pipeline}

The standard RAG architecture consists of two phases: an \textbf{offline
indexing phase} that prepares the knowledge base, and an \textbf{online query
phase} that retrieves and generates.

\begin{figure}[H]
\centering
\begin{tikzpicture}[
    node distance=0.6cm and 1.0cm,
    box/.style={rectangle, draw=deepblue, fill=lightblue, minimum width=2.4cm,
                minimum height=0.65cm, font=\scriptsize\sffamily, rounded corners=2pt,
                thick, align=center},
    indexbox/.style={rectangle, draw=deepblue!60, fill=blue!6, minimum width=2.2cm,
                     minimum height=0.65cm, font=\scriptsize\sffamily, rounded corners=2pt,
                     align=center},
    storebox/.style={rectangle, draw=deepblue!80, fill=deepblue!12, minimum width=2.2cm,
                     minimum height=0.65cm, font=\scriptsize\sffamily, rounded corners=2pt,
                     thick, align=center},
    arr/.style={-{Stealth[length=4pt]}, thick, deepblue!70},
    darr/.style={-{Stealth[length=4pt]}, thick, deepblue!40, dashed},
    lbl/.style={font=\tiny\sffamily, text=deepblue!70},
    phase/.style={font=\scriptsize\sffamily\bfseries, text=gray!60!black}
]

% === OFFLINE INDEXING (top) ===
\node[phase] (offline_lbl) {Offline Indexing};
\node[indexbox, below=0.3cm of offline_lbl] (docs) {Documents};
\node[indexbox, right=0.7cm of docs] (chunk) {Chunking};
\node[indexbox, right=0.7cm of chunk] (embed_idx) {Embedding\\Model};
\node[storebox, right=0.7cm of embed_idx] (vecdb) {Vector DB\\+ Index};

\draw[arr] (docs) -- (chunk);
\draw[arr] (chunk) -- (embed_idx);
\draw[arr] (embed_idx) -- (vecdb);

% === ONLINE QUERY (bottom) ===
\node[phase, below=1.4cm of offline_lbl] (online_lbl) {Online Query};
\node[box, below=0.3cm of online_lbl] (query) {User Query};
\node[box, right=0.7cm of query] (q_enc) {Query\\Encoding};
\node[box, right=0.7cm of q_enc] (retrieve) {Retrieval\\(Top-$K$)};
\node[box, right=0.7cm of retrieve] (rerank) {Re-ranking\\(Top-$k$)};

\node[box, below=0.7cm of rerank] (context) {Context\\Assembly};
\node[box, left=0.7cm of context] (llm) {LLM\\Generation};
\node[box, left=0.7cm of llm, fill=deepblue!15] (answer) {Answer +\\Citations};

\draw[arr] (query) -- (q_enc);
\draw[arr] (q_enc) -- (retrieve);
\draw[arr] (retrieve) -- (rerank);
\draw[arr] (rerank) -- (context);
\draw[arr] (context) -- (llm);
\draw[arr] (llm) -- (answer);

% Connection from vector DB to retrieval
\draw[darr] (vecdb) |- ([yshift=-0.15cm]retrieve.north);

% Separator line
\draw[gray!40, dashed, thick] ([xshift=-0.8cm, yshift=-0.55cm]offline_lbl.south west) --
    ([xshift=0.8cm, yshift=-0.55cm]vecdb.south east);

\end{tikzpicture}
\caption{The complete RAG pipeline.  Documents are chunked, embedded, and
indexed offline.  At query time, the user query is encoded, relevant chunks
are retrieved from the vector database, re-ranked for precision, assembled
into a context window, and fed to the LLM for grounded generation.}
\label{fig:rag-pipeline}
\end{figure}

\paragraph{Offline indexing.}
\begin{enumerate}
    \item \textbf{Document ingestion:} Collect documents from the knowledge
          base (PDFs, HTML, databases, wikis).
    \item \textbf{Chunking:} Split documents into retrieval-sized units
          (Section~\ref{sec:chunking}).
    \item \textbf{Embedding:} Encode each chunk into a dense vector using an
          embedding model.
    \item \textbf{Indexing:} Store embeddings and metadata in a vector database
          with an approximate nearest neighbor (ANN) index.
\end{enumerate}

\paragraph{Online query.}
\begin{enumerate}
    \item \textbf{Query encoding:} Embed the user query using the same (or
          compatible) embedding model.
    \item \textbf{Retrieval:} Find the top-$K$ most similar chunks via ANN
          search (and optionally sparse/hybrid retrieval).
    \item \textbf{Re-ranking:} Apply a more expensive cross-encoder or
          LLM-based re-ranker to select the top-$k$ ($k \ll K$) most
          relevant chunks.
    \item \textbf{Context assembly:} Format the selected chunks into a prompt
          with appropriate instructions.
    \item \textbf{Generation:} The LLM generates a response grounded in the
          retrieved context.
    \item \textbf{Citation:} Optionally, the response includes references to
          specific source chunks.
\end{enumerate}

\begin{keyinsight}[RAG as the Bridge Between Parametric and Non-Parametric Knowledge]
An LLM's parameters encode \emph{parametric knowledge}---facts compressed into
weights during pretraining.  A retrieval system provides \emph{non-parametric
knowledge}---explicit documents that can be updated, verified, and cited.  RAG
combines both: the LLM's language understanding and reasoning capabilities are
grounded by retrieved evidence, producing responses that are more accurate,
more current, and more verifiable than either component alone.  This is why
RAG has become the default architecture for enterprise AI applications---it
gives LLMs access to organizational knowledge without retraining.
\end{keyinsight}


% ============================================================================
% SECTION 2: RETRIEVAL PARADIGMS
% ============================================================================
\section{Retrieval Paradigms}
\label{sec:retrieval-paradigms}

The retrieval component is the backbone of any RAG system.  Choosing the right
retrieval strategy---and understanding the tradeoffs---is one of the most
important design decisions and a frequent interview topic.

\subsection{Sparse Retrieval (BM25)}

Sparse retrieval methods represent queries and documents as high-dimensional
sparse vectors over the vocabulary, where most entries are zero.  The dominant
method remains \textbf{BM25} (Robertson et al., 1994), which extends TF-IDF
with term frequency saturation and document length normalization (see
Chapter~\ref{chap:classical_foundations} for the full derivation).

\paragraph{How it works.}
An \textbf{inverted index} maps each term in the vocabulary to the list of
documents containing that term.  At query time, the system looks up each query
term in the index, retrieves the posting lists, and scores documents using the
BM25 formula.  Only documents sharing at least one term with the query receive
a non-zero score.

\paragraph{Strengths.}
\begin{itemize}
    \item \textbf{Exact lexical matching:} Handles entity names, product IDs,
          technical jargon, and rare terms that dense models often miss.
    \item \textbf{Efficiency:} Inverted index lookup is $O(1)$ per term; no
          GPU required at query time.
    \item \textbf{No training data needed:} Works out of the box on any corpus.
    \item \textbf{Interpretable:} Scores decompose into per-term contributions.
\end{itemize}

\paragraph{Weaknesses.}
\begin{itemize}
    \item \textbf{Vocabulary mismatch:} ``automobile'' and ``car'' share no
          terms, so BM25 scores zero despite identical meaning.
    \item \textbf{No semantic understanding:} Cannot match paraphrases, handle
          synonyms, or bridge conceptual gaps.
    \item \textbf{Bag-of-words assumption:} Ignores word order and context.
\end{itemize}

\subsection{Dense Retrieval (Bi-Encoder)}
\label{sec:dense-retrieval}

Dense retrieval maps queries and documents into a shared dense embedding space
where semantic similarity is measured by vector distance.

\paragraph{Architecture.}
A \textbf{bi-encoder} uses two separate encoders (or a shared encoder) to
independently encode the query and document into fixed-size vectors:
\[
    \vq = E_Q(\text{query}), \qquad \mathbf{d} = E_D(\text{document})
\]
Similarity is computed as cosine similarity or dot product:

\begin{mathresult}{Bi-Encoder Similarity and Contrastive Loss}
\[
    \text{sim}(\vq, \mathbf{d}) = \frac{\vq^\top \mathbf{d}}{\norm{\vq}\,\norm{\mathbf{d}}}
    \qquad \text{or} \qquad
    \text{sim}(\vq, \mathbf{d}) = \vq^\top \mathbf{d}
\]
Training uses the \textbf{InfoNCE contrastive loss} over a batch of queries
$\{q_i\}$ with their positive documents $\{d_i^+\}$:
\[
    \loss = -\frac{1}{N}\sum_{i=1}^{N}
    \log \frac{\exp\!\bigl(\text{sim}(q_i, d_i^+)/\tau\bigr)}
    {\sum_{j=1}^{N} \exp\!\bigl(\text{sim}(q_i, d_j^+)/\tau\bigr)}
\]
where $\tau$ is a temperature hyperparameter and the denominator includes
\textbf{in-batch negatives}---other queries' positives serve as negatives.
\end{mathresult}

\paragraph{Key models.}
\begin{itemize}
    \item \textbf{DPR} (Karpukhin et al., EMNLP 2020): Separate BERT encoders
          for query and passage.  Trained with in-batch negatives plus one BM25
          hard negative per query.  Foundational model for open-domain QA and RAG.
    \item \textbf{E5} (Wang et al., 2022): Instruction-tuned embeddings with
          ``query:'' and ``passage:'' prefixes.  Contrastive pretraining on
          large paired datasets followed by fine-tuning.
    \item \textbf{BGE} (Xiao et al., 2023): Similar instruction-tuned approach
          with strong multilingual support.  Competitive across MTEB benchmarks.
    \item \textbf{GTE} (Li et al., 2023): General Text Embeddings with
          multi-stage training.  Strong out-of-box performance.
    \item \textbf{API-based:} OpenAI Embeddings, Cohere Embed v3, Voyage AI---
          production-ready, require no infrastructure, but introduce API
          dependency and cost.
\end{itemize}

\paragraph{Training details that matter.}
\begin{itemize}
    \item \textbf{Hard negative mining:} In-batch negatives are easy.  Mining
          hard negatives (documents that are similar but not relevant) using
          BM25 or a cross-encoder dramatically improves retrieval quality.
    \item \textbf{Batch size:} Larger batches provide more in-batch negatives,
          improving contrastive learning.  Typical: 512--4096 pairs.
    \item \textbf{Temperature $\tau$:} Lower temperature produces sharper
          discrimination but can cause training instability.  Typical:
          $\tau = 0.01$--$0.1$.
\end{itemize}

\subsection{Learned Sparse Retrieval (SPLADE)}

\textbf{SPLADE} (Formal et al., SIGIR 2021) bridges the gap between sparse and
dense retrieval.  It uses a pretrained language model's MLM head to produce
\emph{learned sparse representations}: for each input token, the MLM head
outputs a distribution over the full vocabulary, effectively performing
\textbf{neural term expansion}.

For example, a document about ``automobile maintenance'' might generate high
weights for terms like ``car,'' ``vehicle,'' ``repair,'' and ``mechanic''---
terms not in the original text but semantically relevant.

\paragraph{Advantages.}
\begin{itemize}
    \item \textbf{Interpretable:} Output is a weighted bag-of-words; you can
          inspect which terms were expanded.
    \item \textbf{Efficient:} Uses inverted index infrastructure (no ANN).
    \item \textbf{Semantic:} Captures synonyms and related terms via expansion.
\end{itemize}

\subsection{Late Interaction (ColBERT)}

\textbf{ColBERT} (Khattab \& Zaharia, SIGIR 2020) introduces a middle ground
between bi-encoders and cross-encoders: \textbf{late interaction}.

\paragraph{How it works.}
Both query and document are independently encoded into \emph{per-token}
embeddings (not a single vector).  At scoring time, each query token's
embedding computes its maximum similarity with any document token embedding,
and these MaxSim scores are summed:
\[
    \text{score}(q, d) = \sum_{i=1}^{|q|}
    \max_{j \in \{1,\ldots,|d|\}} \vq_i^\top \mathbf{d}_j
\]

\paragraph{Tradeoffs.}
ColBERT achieves higher quality than bi-encoders because token-level matching
captures fine-grained semantic alignment.  However, it requires storing
per-token embeddings for every document, resulting in significantly larger
indices (typically $20$--$40\times$ larger than single-vector bi-encoder indices).

\subsection{Cross-Encoder Re-Ranker}

A \textbf{cross-encoder} takes the concatenated (query, document) pair as
input to a single transformer, allowing full bidirectional attention between
query and document tokens:
\[
    \text{score}(q, d) = \text{Linear}\!\bigl(\text{BERT}(\texttt{[CLS]}\; q \;\texttt{[SEP]}\; d)\bigr)
\]

This produces the highest quality scores because the model can perform deep
token-level interaction.  However, the cost is $O(n)$ \emph{per document}:
each (query, document) pair requires a full forward pass.  This makes
cross-encoders impractical for searching a full corpus but ideal for
\emph{re-ranking} a small candidate set (Section~\ref{sec:reranking}).

\subsection{Comparison of Retrieval Methods}

\begin{table}[H]
\centering
\small
\begin{tabularx}{\textwidth}{@{}l l c c c X@{}}
\toprule
\textbf{Method} & \textbf{Mechanism} & \textbf{Quality} & \textbf{Latency} & \textbf{Training?} & \textbf{Key Models} \\
\midrule
BM25 (Sparse)       & Inverted index, lexical  & Moderate & Very fast & None    & Lucene, Elasticsearch \\
Bi-Encoder (Dense)   & Dual encoders, ANN       & High     & Fast      & Required & DPR, E5, BGE, GTE \\
SPLADE (Learned Sparse) & Neural term expansion & High     & Fast      & Required & SPLADE v2, SPLADE++ \\
ColBERT (Late Inter.)  & Token-level MaxSim     & Very high & Moderate & Required & ColBERT v2 \\
Cross-Encoder        & Full attention over pair  & Highest  & Slow      & Required & DeBERTa-v3, BGE re-ranker \\
\bottomrule
\end{tabularx}
\caption{Comparison of retrieval methods.  Quality and latency are relative
within each category.  Cross-encoders are used only for re-ranking, not
full-corpus retrieval.}
\label{tab:retrieval-comparison}
\end{table}

\begin{warningbox}
A common interview mistake is proposing a cross-encoder for full-corpus
retrieval.  Cross-encoders require a forward pass for \emph{every} document in
the corpus, making them $O(N)$ where $N$ is the corpus size.  For a corpus of
10 million documents, this means 10 million forward passes per query---clearly
impractical.  Cross-encoders are used exclusively for re-ranking a small
candidate set (typically 20--100 documents) retrieved by a cheaper method.
\end{warningbox}


% ============================================================================
% SECTION 3: CHUNKING STRATEGIES
% ============================================================================
\section{Chunking Strategies}
\label{sec:chunking}

Retrieval in RAG operates on \textbf{chunks}---sub-document units that are
embedded and retrieved independently.  Chunking is deceptively important:
poorly chunked documents lead to irrelevant retrieval regardless of how good
the embedding model is.

\subsection{Why Chunking Matters}

Documents range from a few sentences to thousands of pages.  Embedding an
entire document into a single vector forces the embedding model to compress
diverse topics into one representation, diluting the signal for any specific
query.  Conversely, embeddings of appropriately-sized chunks can precisely
match query intent.

\subsection{Chunking Methods}

\paragraph{Fixed-size chunking.}
Split text into chunks of a fixed token count (e.g., 512 tokens).  Simple to
implement but may split sentences mid-word or break logical units.  Often
combined with overlap to mitigate boundary effects.

\paragraph{Sentence-based chunking.}
Split at sentence boundaries, grouping consecutive sentences until a token
limit is reached.  Preserves linguistic coherence but sentence boundaries may
not align with topic boundaries.

\paragraph{Recursive (hierarchical) chunking.}
Split first by large structural boundaries (sections, headings), then by
paragraphs, then by sentences if chunks are still too large.  This respects
document structure and is the default in many production frameworks
(e.g., LangChain's \texttt{RecursiveCharacterTextSplitter}).

\paragraph{Semantic chunking.}
Use an embedding model to detect topic shifts: compute sentence embeddings,
identify points where cosine similarity between consecutive sentences drops
below a threshold, and split there.  More expensive but produces topically
coherent chunks.

\paragraph{Document-aware chunking.}
Respect the document's inherent structure: headings, tables, code blocks,
lists.  Avoid splitting a table across two chunks.  Requires document parsing
(e.g., parsing HTML headers, Markdown structure, or PDF layout).

\paragraph{Overlapping chunks.}
Add overlap between consecutive chunks (typically 10--20\% of chunk size)
to prevent information loss at boundaries.  If a relevant sentence falls
at a chunk boundary, overlap ensures it appears in at least one chunk.

\subsection{Chunk Size Tradeoffs}

\begin{table}[H]
\centering
\small
\begin{tabularx}{\textwidth}{@{}l X X@{}}
\toprule
\textbf{Chunk Size} & \textbf{Pros} & \textbf{Cons} \\
\midrule
Small (128--256 tokens) & More precise retrieval; embedding captures specific meaning & May lack sufficient context for generation; more chunks to index \\
Medium (256--512 tokens) & Good balance of precision and context; typical default & May still split some topics \\
Large (512--1024 tokens) & Rich context per chunk; fewer total chunks & Embedding diluted; retrieval less precise \\
\bottomrule
\end{tabularx}
\caption{Chunk size tradeoffs.  The optimal size depends on the domain, query
complexity, and embedding model's maximum input length.}
\label{tab:chunk-size}
\end{table}

\subsection{Parent-Child Retrieval}

A powerful pattern: create small child chunks for precise retrieval, but
return the \textbf{parent chunk} (or full document section) for generation
context.  This gives the LLM the surrounding context it needs to generate
a coherent answer while using the most precise embedding for retrieval.

\begin{warningbox}
Chunk size is the single most undertuned hyperparameter in production RAG
systems.  Many teams default to 512 tokens and never revisit the decision.
In practice, chunk size should be tuned on a held-out evaluation set---
varying chunk size from 128 to 1024 tokens and measuring retrieval recall
can produce 10--20\% quality improvements for zero additional model cost.
Similarly, overlap ratio, chunking method, and the interaction between chunk
size and embedding model should all be treated as hyperparameters.
\end{warningbox}


% ============================================================================
% SECTION 4: VECTOR DATABASES AND INDEXING
% ============================================================================
\section{Vector Databases and Indexing}
\label{sec:vector-db}

Once chunks are embedded into dense vectors, they must be stored and searched
efficiently.  Exact nearest neighbor search is $O(N \cdot d)$ per query (brute
force), which is impractical for millions of chunks.  \textbf{Approximate
Nearest Neighbor} (ANN) algorithms trade a small amount of recall for
orders-of-magnitude speedup.

\subsection{ANN Algorithms}

\paragraph{HNSW (Hierarchical Navigable Small World).}
The dominant ANN algorithm in production vector databases.  HNSW constructs a
multi-layer graph where the top layers contain long-range connections (for
fast navigation) and the bottom layers contain short-range connections (for
precise search).  Search starts at the top layer and greedily navigates toward
the query vector, descending to lower layers for refinement.

\begin{figure}[H]
\centering
\begin{tikzpicture}[
    node distance=0.6cm,
    pt/.style={circle, draw=deepblue, fill=deepblue!20, minimum size=0.3cm,
               inner sep=0pt, font=\tiny\sffamily},
    qpt/.style={circle, draw=deepred, fill=deepred!30, minimum size=0.35cm,
                inner sep=0pt, font=\tiny\sffamily\bfseries},
    edge/.style={draw=deepblue!40, thick},
    searchedge/.style={draw=deepred!70, very thick, -{Stealth[length=3pt]}},
    layer/.style={font=\scriptsize\sffamily\bfseries, text=deepblue!70},
    lbl/.style={font=\tiny\sffamily, text=gray!70}
]

% Layer 2 (top, sparse)
\node[layer] at (-2.2, 3.0) {Layer 2};
\node[pt] (a2) at (0, 3.0) {};
\node[pt] (b2) at (2.5, 3.0) {};
\node[pt] (c2) at (5.0, 3.0) {};
\draw[edge] (a2) -- (b2);
\draw[edge] (b2) -- (c2);

% Layer 1 (mid, moderate)
\node[layer] at (-2.2, 1.5) {Layer 1};
\node[pt] (a1) at (0, 1.5) {};
\node[pt] (b1) at (1.2, 1.5) {};
\node[pt] (c1) at (2.5, 1.5) {};
\node[pt] (d1) at (3.8, 1.5) {};
\node[pt] (e1) at (5.0, 1.5) {};
\draw[edge] (a1) -- (b1);
\draw[edge] (b1) -- (c1);
\draw[edge] (c1) -- (d1);
\draw[edge] (d1) -- (e1);
\draw[edge] (a1) -- (c1);
\draw[edge] (c1) -- (e1);

% Layer 0 (bottom, dense)
\node[layer] at (-2.2, 0) {Layer 0};
\node[pt] (a0) at (0, 0) {};
\node[pt] (b0) at (0.7, 0) {};
\node[pt] (c0) at (1.2, 0) {};
\node[pt] (d0) at (1.9, 0) {};
\node[pt] (e0) at (2.5, 0) {};
\node[pt] (f0) at (3.2, 0) {};
\node[pt] (g0) at (3.8, 0) {};
\node[pt] (h0) at (4.4, 0) {};
\node[pt] (i0) at (5.0, 0) {};
\draw[edge] (a0) -- (b0); \draw[edge] (b0) -- (c0);
\draw[edge] (c0) -- (d0); \draw[edge] (d0) -- (e0);
\draw[edge] (e0) -- (f0); \draw[edge] (f0) -- (g0);
\draw[edge] (g0) -- (h0); \draw[edge] (h0) -- (i0);
\draw[edge] (a0) -- (d0); \draw[edge] (c0) -- (f0);
\draw[edge] (e0) -- (h0); \draw[edge] (f0) -- (i0);

% Query point
\node[qpt] (q) at (3.5, -0.8) {q};
\node[lbl, right=0.1cm of q] {query};

% Search path
\draw[searchedge] (a2) -- (b2);
\draw[searchedge, deepblue!40] (a2) -- (a1);
\draw[searchedge] (b2) -- (c1);
\draw[searchedge, deepblue!40] (c1) -- (d1);
\draw[searchedge] (d1) -- (g0);
\draw[searchedge] (g0) -- (f0);

% Vertical connections (showing hierarchy)
\draw[gray!30, dashed] (a2) -- (a1);
\draw[gray!30, dashed] (b2) -- (c1);
\draw[gray!30, dashed] (c2) -- (e1);
\draw[gray!30, dashed] (a1) -- (a0);
\draw[gray!30, dashed] (c1) -- (e0);
\draw[gray!30, dashed] (e1) -- (i0);
\draw[gray!30, dashed] (b1) -- (c0);
\draw[gray!30, dashed] (d1) -- (g0);

\end{tikzpicture}
\caption{HNSW index structure.  Upper layers have fewer nodes and long-range
connections for coarse navigation; lower layers are dense for precise local
search.  Search starts at the top and descends to the bottom, greedily moving
toward the query.}
\label{fig:hnsw}
\end{figure}

\textbf{Key properties:}
\begin{itemize}
    \item Search complexity: $O(\log N)$ on average.
    \item Build time: $O(N \log N)$.
    \item Memory: stores the full vectors plus the graph structure.
    \item Recall@10 $> 95$\% is typical with proper tuning (ef\_construction,
          M parameters).
    \item Best quality/speed tradeoff for most production workloads.
\end{itemize}

\paragraph{IVF (Inverted File Index).}
Cluster the vectors using k-means into $C$ clusters (Voronoi cells).
At query time, find the $n_{\text{probe}}$ nearest cluster centroids, then
search only within those clusters.  Faster than HNSW for very large corpora
but lower recall at the same latency.

\paragraph{Product Quantization (PQ).}
Compress vectors by splitting each into subvectors and quantizing each
subvector to the nearest centroid in a learned codebook.  Reduces memory by
$4$--$32\times$ but introduces quantization error.  Often combined with IVF
(IVF-PQ) for large-scale systems where memory is the primary constraint.

\subsection{Vector Database Landscape}

\begin{table}[H]
\centering
\small
\begin{tabularx}{\textwidth}{@{}l l l X@{}}
\toprule
\textbf{Database} & \textbf{Type} & \textbf{ANN} & \textbf{Notes} \\
\midrule
Pinecone    & Managed cloud   & Proprietary   & Fully managed, serverless option \\
Weaviate    & Open-source     & HNSW          & Hybrid search (dense + BM25), GraphQL API \\
Milvus      & Open-source     & HNSW, IVF, PQ & Highly scalable, GPU support \\
Qdrant      & Open-source     & HNSW          & Rust-based, fast, rich filtering \\
Chroma      & Open-source     & HNSW          & Lightweight, good for prototyping \\
pgvector    & PostgreSQL ext. & IVF, HNSW     & Familiar SQL interface, co-located with app data \\
FAISS       & Library         & IVF, PQ, HNSW & Facebook's research library; not a database \\
\bottomrule
\end{tabularx}
\caption{Vector database landscape.  Choice depends on scale, operational
requirements, and existing infrastructure.}
\label{tab:vector-dbs}
\end{table}

\subsection{Index Tradeoffs and Metadata Filtering}

Every ANN index involves a three-way tradeoff between \textbf{recall}
(fraction of true nearest neighbors found), \textbf{latency} (query time),
and \textbf{memory} (index size).  Increasing recall requires either more
memory (storing the graph or more centroids) or more latency (exploring more
candidates).

\paragraph{Metadata filtering.}
Production RAG systems often need to filter results by metadata (e.g., date
range, document type, access permissions).  Two strategies exist:

\begin{itemize}
    \item \textbf{Pre-filtering:} Apply metadata filters first to reduce the
          candidate set, then run ANN search on the filtered subset.  Precise
          but may reduce ANN quality if the filtered set is small.
    \item \textbf{Post-filtering:} Run ANN search on the full index, then
          filter results by metadata.  Guarantees ANN quality but may return
          fewer than $k$ results after filtering.
\end{itemize}

Most production databases (Qdrant, Weaviate, Milvus) support efficient
pre-filtering by combining metadata indices with the ANN structure.


% ============================================================================
% SECTION 5: RE-RANKING
% ============================================================================
\section{Re-Ranking}
\label{sec:reranking}

The standard RAG pattern uses a two-stage retrieval pipeline:
\textbf{retrieve many, re-rank to few.}  The first stage (bi-encoder or BM25)
retrieves a broad candidate set (top-100 or top-200) cheaply; the second stage
(cross-encoder or LLM) re-ranks these candidates to select the most relevant
top-$k$ for the LLM context window.

\subsection{Cross-Encoder Re-Rankers}

A cross-encoder (Section~\ref{sec:retrieval-paradigms}) scores each
(query, document) pair with full bidirectional attention.  Re-ranking 100
candidates with a cross-encoder adds $\sim$200--500ms depending on model size
and hardware, but substantially improves retrieval precision.

Popular re-ranker models: BGE-reranker, Cohere Rerank, cross-encoder/ms-marco
models from Sentence-Transformers.

\subsection{LLM-Based Re-Ranking}

Increasingly, LLMs themselves are used as re-rankers.  Three paradigms exist:

\begin{itemize}
    \item \textbf{Pointwise:} Score each document independently (``Rate the
          relevance of this document to the query on a scale of 1--10'').
    \item \textbf{Pairwise:} Compare pairs of documents (``Which document is
          more relevant to the query?'').
    \item \textbf{Listwise:} Present the full candidate list and ask the LLM
          to rank them (``Rank these documents by relevance'').
\end{itemize}

LLM re-ranking is expensive but can achieve very high quality.  It is most
useful when the candidate set is small (10--20 documents) or when the domain
requires complex relevance judgments that cross-encoders struggle with.

\subsection{Hybrid Retrieval and Reciprocal Rank Fusion}

The most robust production RAG systems combine multiple retrieval methods:
typically BM25 (sparse) + bi-encoder (dense), then fuse the rankings.
\textbf{Reciprocal Rank Fusion} (RRF) is the standard fusion method:

\begin{mathresult}{Reciprocal Rank Fusion (RRF)}
Given $m$ rankers, each producing a ranking over documents, the fused score for
document $d$ is:
\[
    \text{RRF}(d) = \sum_{i=1}^{m} \frac{1}{k + \text{rank}_i(d)}
\]
where $k$ is a constant (typically $k = 60$) that mitigates the impact of high
rankings from any single ranker, and $\text{rank}_i(d)$ is the rank of document
$d$ in the $i$-th ranker's output.
\end{mathresult}

\paragraph{Why hybrid works.}
BM25 excels at exact lexical matching (entity names, technical terms); dense
retrieval excels at semantic matching (paraphrases, conceptual similarity).
Their failure modes are largely complementary.  Hybrid retrieval with RRF
consistently outperforms either method alone across diverse benchmarks (BEIR,
MTEB).

\subsection{Distillation: Best of Both Worlds}

An advanced technique: use a high-quality cross-encoder to score a large set
of (query, document) pairs, then train a bi-encoder to match these scores.
The resulting bi-encoder inherits much of the cross-encoder's quality while
retaining the bi-encoder's efficiency for online retrieval.  This is a common
pattern at companies that can afford the offline computation.


% ============================================================================
% SECTION 6: ADVANCED RAG PATTERNS
% ============================================================================
\section{Advanced RAG Patterns}
\label{sec:advanced-rag}

The basic retrieve-then-generate pattern has well-known failure modes:
ambiguous queries, multi-hop reasoning requirements, and unnecessary retrieval
for simple questions.  Advanced RAG patterns address these limitations.

\subsection{Query Transformation}

\paragraph{Query expansion.}
Add related terms or rephrasings to the original query before retrieval.
Simple approaches use a thesaurus or WordNet; modern approaches use an LLM
to generate expanded queries.  This increases recall by casting a wider
retrieval net.

\paragraph{Query decomposition.}
For complex multi-part questions (``Compare the revenue growth of Apple and
Google in Q3 2024''), decompose into sub-queries (``Apple revenue Q3 2024,''
``Google revenue Q3 2024''), retrieve for each sub-query separately, then
synthesize.

\paragraph{HyDE (Hypothetical Document Embeddings).}
Gao et al.\ (2022) proposed a counterintuitive but effective strategy: instead
of encoding the \emph{query}, generate a \textbf{hypothetical answer} using the
LLM (without retrieval), then encode that hypothetical answer as the retrieval
query.

\begin{figure}[H]
\centering
\begin{tikzpicture}[
    node distance=0.5cm and 0.8cm,
    box/.style={rectangle, draw=deepblue, fill=lightblue, minimum width=2.2cm,
                minimum height=0.6cm, font=\scriptsize\sffamily, rounded corners=2pt,
                thick, align=center},
    arr/.style={-{Stealth[length=4pt]}, thick, deepblue!70},
    lbl/.style={font=\tiny\sffamily, text=deepblue!70}
]

\node[box] (query) {User Query};
\node[box, right=0.8cm of query] (llm1) {LLM\\(no retrieval)};
\node[box, right=0.8cm of llm1] (hypo) {Hypothetical\\Answer};
\node[box, right=0.8cm of hypo] (embed) {Embedding\\Model};
\node[box, below=0.7cm of embed] (retrieve) {Vector\\Retrieval};
\node[box, left=0.8cm of retrieve] (realdocs) {Real\\Documents};
\node[box, left=0.8cm of realdocs] (llm2) {LLM\\(grounded)};
\node[box, left=0.8cm of llm2, fill=deepblue!15] (ans) {Final\\Answer};

\draw[arr] (query) -- (llm1);
\draw[arr] (llm1) -- (hypo);
\draw[arr] (hypo) -- (embed);
\draw[arr] (embed) -- (retrieve);
\draw[arr] (retrieve) -- (realdocs);
\draw[arr] (realdocs) -- (llm2);
\draw[arr] (llm2) -- (ans);

\node[lbl, above=0.05cm of llm1] {generate};
\node[lbl, above=0.05cm of embed] {encode};
\node[lbl, right=0.05cm of retrieve] {ANN search};

\end{tikzpicture}
\caption{HyDE pipeline.  The LLM generates a hypothetical answer (which may be
factually wrong), but its embedding is closer to relevant documents in embedding
space than the original short query.  Real documents are then retrieved and used
to generate a grounded final answer.}
\label{fig:hyde}
\end{figure}

\textbf{Why HyDE works:}  Short queries (``What is RLHF?'') are
poorly represented in embedding space because they lack the vocabulary and
structure of passage-length text.  A hypothetical answer, even if inaccurate,
is a document-length text that uses the same vocabulary, style, and structure as
real documents in the corpus, making it a better retrieval probe.

\textbf{When HyDE fails:}  For queries about specific entities, numbers, or
facts where the LLM's hallucinated answer might be in a completely different
semantic region than the true answer.

\subsection{Multi-Hop Retrieval}

Some questions cannot be answered from a single retrieval step.  For example,
``Who directed the film that won the Palme d'Or in 2024?'' requires:
(1)~retrieve the 2024 Palme d'Or winner, (2)~extract the film name,
(3)~retrieve information about that film's director.

Multi-hop RAG performs iterative retrieval: the LLM reads the first set of
retrieved documents, identifies what additional information is needed,
formulates a new query, retrieves again, and repeats until it has sufficient
information.

\subsection{Self-RAG: When Not to Retrieve}

Asai et al.\ (2023) proposed \textbf{Self-RAG}, where the model learns to
decide \emph{whether} retrieval is needed.  The model generates special
reflection tokens that indicate: (a)~whether retrieval is necessary for this
query, (b)~whether retrieved passages are relevant, and (c)~whether the
generated response is supported by the evidence.

This avoids unnecessary retrieval for simple factual questions the model
already knows (``What is the capital of France?'') while triggering retrieval
for knowledge-intensive or time-sensitive queries.

\subsection{Agentic RAG}

The most flexible RAG pattern treats retrieval as a \textbf{tool} within an
agentic framework.  The LLM can:
\begin{itemize}
    \item Decide whether to retrieve at all.
    \item Choose which knowledge base or data source to query.
    \item Formulate and reformulate queries based on intermediate results.
    \item Chain multiple retrieval steps with reasoning between them.
    \item Use other tools (calculators, code execution, APIs) alongside
          retrieval.
\end{itemize}

This is the paradigm used by systems like Perplexity, ChatGPT with browsing,
and enterprise copilots that need to reason over multiple data sources.

\subsection{Context Compression}

When many chunks are retrieved, they may collectively exceed the LLM's context
window or contain redundant information.  \textbf{Context compression}
summarizes or filters retrieved chunks before feeding them to the LLM:
\begin{itemize}
    \item \textbf{Extractive compression:} Select only the most relevant
          sentences from each chunk (using a smaller model as a filter).
    \item \textbf{Abstractive compression:} Summarize retrieved chunks into a
          concise context using a fast LLM.
    \item \textbf{LLMLingua / AutoCompressor:} Learned token-level compression
          that reduces context length while preserving information.
\end{itemize}

\subsection{Citation and Attribution}

Production RAG systems must support \textbf{citation}: the model's response
should reference specific source chunks so users can verify claims.  Strategies
include:
\begin{itemize}
    \item \textbf{Inline citation:} The model generates bracketed references
          (e.g., ``Revenue grew 15\% [Source 3]'') during generation.
    \item \textbf{Post-hoc attribution:} After generation, a separate model
          maps each claim in the response to supporting chunks.
    \item \textbf{Prompt-based:} Instruct the LLM to cite sources in the system
          prompt (``Always cite the source number for each claim'').
\end{itemize}


% ============================================================================
% SECTION 7: RAG EVALUATION
% ============================================================================
\section{RAG Evaluation}
\label{sec:rag-evaluation}

Evaluating a RAG system is a \textbf{three-layer problem}: you must separately
assess retrieval quality, generation quality, and end-to-end faithfulness.
Many teams make the mistake of evaluating only the final output without
diagnosing whether failures originate in retrieval or generation.

\subsection{Retrieval Metrics}

These metrics evaluate whether the retrieval component finds relevant documents.

\begin{table}[H]
\centering
\small
\begin{tabularx}{\textwidth}{@{}l X l@{}}
\toprule
\textbf{Metric} & \textbf{What It Measures} & \textbf{Use When} \\
\midrule
Recall@$K$ & Fraction of relevant docs in top $K$ & Candidate generation \\
Precision@$K$ & Fraction of top $K$ that are relevant & Small $K$ (top-3, top-5) \\
MRR & $1/\text{rank}$ of first relevant result & Single-answer queries \\
MAP & Average precision across recall levels & Ranked list quality \\
NDCG@$K$ & Graded relevance with position discount & Graded relevance levels \\
Hit Rate@$K$ & Did any relevant doc appear in top $K$? & Binary relevance \\
\bottomrule
\end{tabularx}
\caption{Retrieval evaluation metrics.  For RAG, Recall@$K$ is typically the
most important---if the relevant chunk is not retrieved, the LLM cannot use it.}
\label{tab:retrieval-metrics}
\end{table}

\subsection{Generation Metrics}

Given retrieved context, does the LLM generate a good response?

\begin{itemize}
    \item \textbf{Factual accuracy:} Does the response contain correct facts?
          Measured by human evaluation or automated fact-checking against a
          knowledge base.
    \item \textbf{Relevance:} Does the response address the user's question?
          Measured by LLM-as-judge or human rating.
    \item \textbf{Completeness:} Does the response cover all aspects of the
          question?  Particularly important for multi-part questions.
    \item \textbf{Fluency and coherence:} Is the response well-written?
          Usually high for modern LLMs; more relevant for smaller models.
\end{itemize}

\subsection{Faithfulness: The Critical Metric}

The most important RAG-specific metric is \textbf{faithfulness}: does the
response use \emph{only} information from the retrieved context, without
adding hallucinated claims?

\begin{itemize}
    \item \textbf{NLI-based:} Use a Natural Language Inference model to check
          whether each claim in the response is entailed by the retrieved
          context.
    \item \textbf{LLM-as-judge:} Prompt an LLM to verify each claim against
          the context (``Does the context support this claim? Yes/No'').
    \item \textbf{Citation verification:} Check whether cited sources actually
          contain the information attributed to them.
\end{itemize}

\subsection{End-to-End Frameworks}

\textbf{RAGAS} (Retrieval Augmented Generation Assessment) provides a
framework for end-to-end RAG evaluation with four metrics:
\begin{enumerate}
    \item \textbf{Answer relevancy:} Is the answer relevant to the question?
    \item \textbf{Faithfulness:} Is the answer grounded in the context?
    \item \textbf{Context precision:} Are the retrieved contexts relevant?
    \item \textbf{Context recall:} Do the retrieved contexts cover the answer?
\end{enumerate}

\begin{keyinsight}[RAG Evaluation Is a Three-Layer Problem]
When a RAG system produces a bad answer, the root cause can be at any of three
layers: (1)~\textbf{Retrieval failure}---the relevant document was not
retrieved (fix: better embedding model, hybrid retrieval, query transformation).
(2)~\textbf{Ranking failure}---the relevant document was retrieved but ranked
too low to appear in the context (fix: re-ranking, larger $K$).
(3)~\textbf{Generation failure}---the relevant document was in the context but
the LLM ignored it, misinterpreted it, or hallucinated anyway (fix: prompt
engineering, model upgrade, instruction tuning for grounding).
Diagnosing which layer failed is the first step to improving a RAG system,
and interviewers expect you to articulate this decomposition.
\end{keyinsight}


% ============================================================================
% SECTION 8: RAG VS FINE-TUNING VS LONG CONTEXT
% ============================================================================
\section{RAG vs.\ Fine-Tuning vs.\ Long Context}
\label{sec:rag-vs-ft}

A common interview question asks when to use RAG versus fine-tuning versus
simply stuffing all documents into a long context window.  These are not
mutually exclusive---they can be combined---but understanding the tradeoffs
is critical for system design.

\subsection{Decision Framework}

\begin{table}[H]
\centering
\small
\begin{tabularx}{\textwidth}{@{}l X X X@{}}
\toprule
\textbf{Dimension} & \textbf{RAG} & \textbf{Fine-Tuning} & \textbf{Long Context} \\
\midrule
\textbf{Knowledge source} & External docs, databases & Encoded in model weights & Documents in prompt \\
\textbf{Knowledge update} & Index update (minutes) & Retraining (hours--days) & Upload at query time \\
\textbf{Attribution} & Natural (cite chunks) & Impossible & Possible (cite passages) \\
\textbf{Cost per query} & Retrieval + generation & Generation only & Generation (higher token cost) \\
\textbf{Corpus size} & Unlimited & Limited by training & Limited by context window \\
\textbf{Latency} & Retrieval adds 50--200ms & Lowest & Higher (long prompt) \\
\textbf{Best for} & Dynamic knowledge, compliance & Behavior, style, format & Small corpora, simple QA \\
\textbf{Failure mode} & Retrieval misses & Knowledge not learned & Lost in the middle \\
\bottomrule
\end{tabularx}
\caption{Decision framework for RAG vs.\ fine-tuning vs.\ long context.
These approaches can be combined: fine-tune for domain adaptation + RAG for
knowledge grounding.}
\label{tab:rag-vs-ft}
\end{table}

\subsection{When to Use Each}

\paragraph{Use RAG when:}
\begin{itemize}
    \item Knowledge changes frequently (news, product catalogs, documentation).
    \item Attribution and verifiability are requirements (legal, medical, compliance).
    \item The knowledge base is large (millions of documents).
    \item You need to ground the model in specific organizational knowledge.
\end{itemize}

\paragraph{Use fine-tuning when:}
\begin{itemize}
    \item You need to change the model's \emph{behavior}, not its knowledge
          (tone, format, following domain-specific instructions).
    \item Consistent stylistic adaptation is required (writing like a specific
          brand voice).
    \item The domain has specialized terminology or reasoning patterns
          (biomedical, legal, financial).
    \item You want to distill a larger model into a smaller, faster one.
\end{itemize}

\paragraph{Use long context when:}
\begin{itemize}
    \item The corpus is small enough to fit in a context window (a single
          document, a handful of pages).
    \item Real-time analysis is needed (``summarize this document'').
    \item Pipeline simplicity is paramount (no retrieval infrastructure).
    \item The task requires holistic understanding of the full document.
\end{itemize}

\paragraph{Combine them:}
\begin{itemize}
    \item \textbf{Fine-tune + RAG:} Fine-tune for domain style and instruction
          following; RAG for knowledge grounding.  This is the most common
          production pattern.
    \item \textbf{Long context + RAG:} Retrieve a large set of documents and
          use a long-context model to process all of them.  Reduces the
          ``lost in the middle'' risk by retrieving more and letting the model
          attend to everything.
\end{itemize}

\begin{keyinsight}[The ``Lost in the Middle'' Problem]
Liu et al.\ (2023) demonstrated that LLMs disproportionately attend to
information at the \emph{beginning} and \emph{end} of the context window,
with degraded recall for information in the middle.  This has direct
implications for RAG: the order in which retrieved chunks are placed in the
prompt matters.  Mitigation strategies include: (1)~placing the most relevant
chunks at the beginning and end of the context; (2)~using fewer, more
precisely retrieved chunks rather than many; (3)~using models specifically
trained for long-context faithfulness; (4)~iterating with multiple retrieval
rounds to validate key facts.
\end{keyinsight}


% ============================================================================
% SECTION 9: INTERVIEW QUESTIONS
% ============================================================================
\section{Interview Questions}
\label{sec:rag-questions}

%% QUESTION 1 %%
\begin{interviewq}
\textbf{Q: Design a RAG system for a customer support chatbot that answers
questions using your company's knowledge base. Walk through the key components
and design decisions.}

\badgeLfive{L5} \quad \badgetype{System Design} \quad \badgefreq{\freqhigh}

\stronganswer
The system has two phases.  \textbf{Offline indexing:} ingest knowledge base
articles, chunk them using recursive chunking (respect headings and sections,
target 256--512 tokens with 10\% overlap), embed each chunk using an
instruction-tuned model like E5 or BGE, and index embeddings in a vector
database (Qdrant or pgvector if we want colocation with metadata).  Store
metadata: article ID, section, last-updated date, product category.

\textbf{Online query:} encode the customer query with the same embedding model,
perform hybrid retrieval (BM25 for exact product names/error codes + dense
retrieval for semantic matching), fuse with RRF, re-rank the top-50 candidates
with a cross-encoder to select top-5, assemble a prompt with system instructions
(``Answer using only the provided context; cite source articles; say `I don't
know' if the context is insufficient''), and generate with the LLM.

\textbf{Key design decisions:} (1)~Hybrid retrieval is critical---customer
queries often contain exact error codes (BM25 wins) mixed with natural language
descriptions (dense wins).  (2)~Metadata filtering by product category
dramatically improves precision.  (3)~A fallback to human agents when
confidence is low.  (4)~Evaluation: track retrieval recall on a labeled QA set,
faithfulness via NLI checks, and user satisfaction ratings.  (5)~Monitoring:
log queries where the model says ``I don't know'' to identify knowledge gaps.

\redflags
\begin{itemize}
    \item Proposes only dense or only sparse retrieval without considering hybrid
    \item Ignores chunking strategy or uses a single chunk size without justification
    \item No mention of evaluation, monitoring, or fallback mechanisms
    \item Proposes a cross-encoder for initial retrieval (does not scale)
\end{itemize}

\interviewtest{Whether the candidate can design a complete end-to-end RAG
system with production awareness: data pipeline, retrieval strategy,
generation, evaluation, and monitoring---not just ``use an embedding model
and an LLM.''}

\goodgreat
\begin{itemize}
    \item \badgeLfive{L5} Describes the full pipeline with reasonable component choices,
          mentions hybrid retrieval and re-ranking
    \item \badgeLsix{L6} Discusses metadata filtering, chunk size tuning, the
          offline/online split, evaluation strategy, and failure modes
    \item \badgeLseven{L7} Designs for multi-turn conversation (session context),
          discusses guardrails for hallucination, proposes an active learning loop
          to improve retrieval from user feedback, and addresses scaling (sharding
          the index, caching frequent queries)
\end{itemize}

\followups{How would you handle multi-turn conversations where context from
previous turns matters?  How would you detect and mitigate hallucination in
real time?  How would you handle knowledge base updates without downtime?}
\end{interviewq}

\bigskip

%% QUESTION 2 %%
\begin{interviewq}
\textbf{Q: Compare BM25 vs.\ dense retrieval. When would you use each? When
would you combine them?}

\badgeLfive{L5} \quad \badgetype{Trade-off} \quad \badgefreq{\freqhigh}

\stronganswer
\textbf{BM25} uses lexical matching via an inverted index: it scores documents
based on query term overlap, weighted by TF-IDF-like statistics.  Strengths:
handles exact terms (entity names, product IDs, error codes), no training data
needed, extremely fast, interpretable.  Weakness: no semantic understanding---
fails on vocabulary mismatch (``automobile'' vs.\ ``car'').

\textbf{Dense retrieval} (bi-encoder) embeds queries and documents in a shared
vector space using trained neural encoders.  Strengths: captures semantic
similarity, handles paraphrases and conceptual queries.  Weakness: may miss
exact lexical matches, requires training data and GPU infrastructure, less
interpretable.

\textbf{When to use BM25:} low-resource settings (no training data), domains
with precise terminology (medical codes, legal statutes), when exact matching
is critical.  \textbf{When to use dense:} semantic search, FAQ matching, when
queries and documents use different vocabulary.  \textbf{Combine them:} almost
always in production.  Hybrid retrieval with RRF consistently outperforms
either alone because their failure modes are complementary.  BM25 catches
what dense misses (rare terms) and vice versa.

\redflags
\begin{itemize}
    \item Says BM25 is ``obsolete'' or ``only for legacy systems''
    \item Cannot explain the vocabulary mismatch problem
    \item Does not mention the option to combine them
    \item Cannot articulate when BM25 would outperform dense retrieval
\end{itemize}

\interviewtest{Whether the candidate understands that retrieval is not
one-size-fits-all and can reason about when simpler methods outperform neural
ones---a sign of production maturity, not just knowledge of the latest models.}

\goodgreat
\begin{itemize}
    \item \badgeLfive{L5} Clearly articulates the strengths and weaknesses of each,
          recommends hybrid with RRF
    \item \badgeLsix{L6} Discusses SPLADE as a learned sparse alternative, mentions
          ColBERT for higher quality, and quantifies the tradeoff (BM25 wins by
          5--10 points on entity-heavy queries, dense wins on paraphrase-heavy queries)
    \item \badgeLseven{L7} Discusses the BEIR benchmark (out-of-domain generalization),
          notes that BM25 is surprisingly competitive on many domains, and proposes
          adaptive retrieval strategies that route queries to BM25 vs.\ dense based
          on query characteristics
\end{itemize}

\followups{What is SPLADE and how does it bridge sparse and dense?  How would
you decide the RRF constant $k$?  What is the BEIR benchmark and why does it
matter for retrieval?}
\end{interviewq}

\bigskip

%% QUESTION 3 %%
\begin{interviewq}
\textbf{Q: What chunking strategies exist for RAG? How do you choose chunk
size?}

\badgeLfive{L5} \quad \badgetype{System Design} \quad \badgefreq{\freqhigh}

\stronganswer
Main strategies: (1)~\textbf{Fixed-size}: split at $N$ tokens with overlap;
simple but may break sentences or logical units.  (2)~\textbf{Sentence-based}:
group sentences to a token limit; preserves linguistic coherence.
(3)~\textbf{Recursive/hierarchical}: split by headings, then paragraphs, then
sentences if needed; respects document structure.  (4)~\textbf{Semantic}:
detect topic shifts via embedding similarity drops; produces topically coherent
chunks.  (5)~\textbf{Document-aware}: respect tables, code blocks, lists as
atomic units.

\textbf{Choosing chunk size:} smaller chunks (128--256 tokens) give more
precise retrieval but may lack context for generation; larger chunks (512--1024)
provide richer context but dilute the embedding signal.  The sweet spot is
typically 256--512 tokens with 10--20\% overlap, but this must be tuned
empirically.  The \textbf{parent-child pattern} solves the precision-context
tradeoff: retrieve small child chunks for precision, return the parent section
for context.  Chunk size should be treated as a hyperparameter tuned on a
held-out retrieval evaluation set.

\redflags
\begin{itemize}
    \item Only knows fixed-size chunking
    \item Cannot explain why chunk size matters for retrieval quality
    \item Does not mention overlap or the parent-child pattern
    \item Treats chunking as a solved problem rather than a hyperparameter
\end{itemize}

\interviewtest{Whether the candidate understands the practical engineering of
RAG systems beyond the model level.  Chunking is ``unglamorous'' but has
enormous impact on system quality.}

\goodgreat
\begin{itemize}
    \item \badgeLfive{L5} Describes 3+ chunking strategies and explains the size tradeoff
    \item \badgeLsix{L6} Mentions parent-child retrieval, discusses interaction between
          chunk size and embedding model's max input length, and proposes evaluation-driven
          tuning
    \item \badgeLseven{L7} Discusses multi-granularity indexing (chunks at multiple sizes),
          late chunking (encode full document then extract chunk embeddings), and
          agentic approaches where the system dynamically adjusts retrieval granularity
\end{itemize}

\followups{How does chunk size interact with the embedding model's context
length?  What is late chunking?  How would you handle chunking for
tables and images?}
\end{interviewq}

\bigskip

%% QUESTION 4 %%
\begin{interviewq}
\textbf{Q: How would you evaluate a RAG system end-to-end? What metrics would
you track?}

\badgeLfive{L5} \quad \badgetype{Evaluation} \quad \badgefreq{\freqhigh}

\stronganswer
RAG evaluation must be decomposed into three layers.  \textbf{Layer 1:
Retrieval quality.}  Metrics: Recall@$K$ (most important---if the relevant
chunk is not retrieved, the LLM cannot use it), MRR, NDCG@$K$.  Requires a
labeled set of (query, relevant chunk) pairs.

\textbf{Layer 2: Generation quality.}  Given retrieved context, does the model
produce a good answer?  Metrics: answer relevancy, completeness, and coherence.
Measured via human evaluation or LLM-as-judge.

\textbf{Layer 3: Faithfulness.}  The RAG-specific critical metric: does the
answer contain \emph{only} information supported by the retrieved context?
Measured via NLI-based verification, LLM-as-judge faithfulness scoring, or
citation verification.  Frameworks like RAGAS compute these metrics
automatically.

\textbf{End-to-end:} answer correctness (is the final answer factually
correct?) and user satisfaction (thumbs up/down, CSAT scores).

Key insight: when the system produces a wrong answer, you must diagnose
\emph{which layer} failed.  Retrieval failure needs better embeddings or hybrid
search; ranking failure needs re-ranking; generation failure needs better
prompting or model.

\redflags
\begin{itemize}
    \item Evaluates only the final answer without decomposing into layers
    \item Does not mention faithfulness as a metric
    \item Cannot name specific metrics (Recall@$K$, NDCG, MRR)
    \item ``We'll just use human evaluation'' without a concrete plan
\end{itemize}

\interviewtest{Whether the candidate has a systematic approach to evaluation
that enables debugging.  Production RAG teams spend more time on evaluation
than on model selection.}

\goodgreat
\begin{itemize}
    \item \badgeLfive{L5} Decomposes evaluation into retrieval, generation, and
          faithfulness layers with specific metrics for each
    \item \badgeLsix{L6} Discusses LLM-as-judge tradeoffs (verbosity bias, position
          bias), mentions RAGAS, and proposes a labeled evaluation set strategy
    \item \badgeLseven{L7} Designs a continuous evaluation pipeline with automated
          regression detection, discusses calibrating LLM-as-judge against human
          annotations, and proposes adversarial evaluation (queries designed to
          trigger hallucination)
\end{itemize}

\followups{How would you build a labeled evaluation set for RAG?  What are
the biases of LLM-as-judge evaluation?  How do you detect quality regression
in production?}
\end{interviewq}

\bigskip

%% QUESTION 5 %%
\begin{interviewq}
\textbf{Q: What is HyDE (Hypothetical Document Embeddings)? When would it help
and when would it fail?}

\badgeLsix{L6} \quad \badgetype{Conceptual} \quad \badgefreq{\freqmed}

\stronganswer
HyDE uses the LLM to generate a \emph{hypothetical answer} to the query
(without retrieval), then embeds the hypothetical answer as the retrieval
query instead of the original question.  The intuition: short queries are
poor retrieval probes because they lack the vocabulary and structure of the
documents in the index.  A hypothetical answer, even if factually incorrect,
is a passage-length text that uses similar vocabulary and structure as real
passages, making it a better embedding for dense retrieval.

\textbf{When it helps:} abstract or conceptual queries (``How does quantum
entanglement work?''), queries where the user's phrasing differs significantly
from document language, and domains where queries are short but documents are
long.

\textbf{When it fails:} (1)~Factual queries about specific entities or numbers,
where the LLM's hallucinated answer might be in a completely wrong semantic
region.  (2)~When the LLM has very poor knowledge of the domain, so the
hypothetical answer is semantically far from any real document.  (3)~Adds
latency (one extra LLM call per query).  (4)~Queries that are already
document-like (copy-pasted error messages, long-form questions).

\redflags
\begin{itemize}
    \item Cannot explain why embedding the hypothetical answer helps retrieval
    \item Thinks HyDE eliminates hallucination (it only improves retrieval;
          the final answer still comes from real documents)
    \item Does not mention failure cases
\end{itemize}

\interviewtest{Whether the candidate can reason about non-obvious RAG patterns
and critically evaluate when they apply.  HyDE is a creative solution that
tests understanding of embedding spaces.}

\goodgreat
\begin{itemize}
    \item \badgeLfive{L5} Describes HyDE correctly and mentions the query-document
          asymmetry problem
    \item \badgeLsix{L6} Explains the embedding space intuition (document-length
          texts cluster near similar documents), discusses failure modes, and
          mentions the latency tradeoff
    \item \badgeLseven{L7} Connects HyDE to query expansion as a general strategy,
          discusses multi-hypothesis retrieval (generate multiple hypothetical
          answers, retrieve for each, merge), and compares to query decomposition
          approaches
\end{itemize}

\followups{How does HyDE relate to query expansion?  Could you generate
multiple hypothetical answers and merge the retrieval results?  How would
you decide at query time whether to use HyDE or direct retrieval?}
\end{interviewq}

\bigskip

%% QUESTION 6 %%
\begin{interviewq}
\textbf{Q: How do you handle the ``lost in the middle'' problem in RAG?}

\badgeLsix{L6} \quad \badgetype{System Design} \quad \badgefreq{\freqmed}

\stronganswer
Liu et al.\ (2023) showed that LLMs perform best on information at the
\emph{beginning} and \emph{end} of the context window, with significant
degradation for information in the middle---a U-shaped recall curve.  This
directly impacts RAG because the placement of retrieved chunks in the prompt
affects whether the model actually uses them.

\textbf{Mitigation strategies:}
(1)~\textbf{Position-aware ordering:} Place the most relevant chunks at the
beginning and end of the context, less relevant ones in the middle.
(2)~\textbf{Retrieve fewer, better chunks:} Instead of stuffing 20 chunks,
use aggressive re-ranking to select the top 3--5 most relevant, reducing the
risk of important information being lost.  (3)~\textbf{Context compression:}
Summarize retrieved chunks before insertion to reduce total length.
(4)~\textbf{Multi-pass generation:} Have the model process chunks in separate
passes and aggregate the results.  (5)~\textbf{Use models trained for
long-context faithfulness:} Some models (Claude, Gemini 1.5) are specifically
trained to attend uniformly across long contexts.
(6)~\textbf{Iterative retrieval:} Instead of one large retrieval, perform
multiple smaller retrievals with intermediate reasoning.

\redflags
\begin{itemize}
    \item Has not heard of the ``lost in the middle'' phenomenon
    \item Proposes only ``use a longer context window'' without addressing the
          attention distribution problem
    \item Cannot articulate specific mitigation strategies
\end{itemize}

\interviewtest{Whether the candidate understands the practical limitations of
long-context LLMs and can design around them.  This tests production RAG
awareness beyond the textbook architecture.}

\goodgreat
\begin{itemize}
    \item \badgeLfive{L5} Describes the problem and 2--3 mitigation strategies
    \item \badgeLsix{L6} References the Liu et al.\ paper, discusses the U-shaped
          recall curve, and proposes evaluation-driven chunk ordering
    \item \badgeLseven{L7} Discusses the relationship to attention mechanisms
          (why the U-shape exists), evaluates mitigation strategies quantitatively,
          and proposes architectures that address this at the model level (e.g.,
          structured attention, landmark tokens)
\end{itemize}

\followups{Why does the U-shaped recall pattern exist mechanistically?  How
would you evaluate whether your mitigation is working?  Does this problem
improve with newer models?}
\end{interviewq}

\bigskip

%% QUESTION 7 %%
\begin{interviewq}
\textbf{Q: When should you use RAG vs.\ fine-tuning vs.\ long context?}

\badgeLfive{L5} \quad \badgetype{Trade-off} \quad \badgefreq{\freqhigh}

\stronganswer
These serve different purposes and are not mutually exclusive.

\textbf{RAG:} when knowledge changes frequently, when you need attribution and
citations, when the knowledge base is too large for any context window, and when
verifiability matters (legal, medical, compliance).  RAG's knowledge can be
updated by re-indexing without retraining.

\textbf{Fine-tuning:} when you need to change the model's \emph{behavior}---
its tone, format, instruction-following style, or domain reasoning patterns.
Fine-tuning encodes knowledge into weights, so it cannot be easily updated
and provides no attribution.  Best for: style adaptation, domain-specific
instruction following, distillation.

\textbf{Long context:} when the corpus is small enough to fit in a context
window, when you need real-time document analysis (``summarize this PDF''),
or when pipeline simplicity is paramount.  Limitations: cost scales linearly
with context length, ``lost in the middle'' problem, context window is still
finite.

\textbf{Combine:} fine-tune + RAG is the most common production pattern
(fine-tune for domain behavior, RAG for knowledge grounding).  Long context
+ RAG is emerging: retrieve a larger set and let the model process everything.

\redflags
\begin{itemize}
    \item Treats these as mutually exclusive
    \item Says ``just use RAG for everything'' or ``just fine-tune''
    \item Cannot articulate when fine-tuning beats RAG (behavior vs.\ knowledge)
    \item Ignores cost, latency, and updateability considerations
\end{itemize}

\interviewtest{Whether the candidate can make nuanced engineering decisions
about the right tool for the job.  This is fundamentally about matching the
solution to the problem constraints.}

\goodgreat
\begin{itemize}
    \item \badgeLfive{L5} Clearly articulates when to use each and mentions combining them
    \item \badgeLsix{L6} Discusses the knowledge vs.\ behavior distinction, provides
          concrete examples of each regime, and mentions cost/latency tradeoffs
    \item \badgeLseven{L7} Proposes a decision tree based on problem characteristics
          (knowledge freshness, corpus size, attribution requirements, latency budget),
          discusses how the tradeoffs shift as context windows grow, and mentions
          techniques like RAFT (Retrieval Augmented Fine-Tuning) that blur the boundary
\end{itemize}

\followups{What is RAFT (Retrieval Augmented Fine-Tuning)?  How do the
tradeoffs change as context windows grow to 1M+ tokens?  When would you
fine-tune the retrieval model vs.\ the generation model?}
\end{interviewq}

\bigskip

%% QUESTION 8 %%
\begin{interviewq}
\textbf{Q: Explain the bi-encoder vs.\ cross-encoder tradeoff in retrieval.
How would you use both in a production system?}

\badgeLfive{L5} \quad \badgetype{Conceptual} \quad \badgefreq{\freqhigh}

\stronganswer
A \textbf{bi-encoder} independently encodes query and document into vectors,
then computes similarity via dot product or cosine.  Because documents can be
pre-encoded offline and stored in an ANN index, retrieval is $O(\log N)$ over
the corpus.  However, no cross-attention means the model cannot perform deep
query-document interaction.

A \textbf{cross-encoder} concatenates (query, document) and processes them
jointly through a full transformer with bidirectional attention.  This allows
rich token-level interaction (the model can see both query and document tokens
simultaneously), producing much higher quality relevance scores.  However,
it requires a forward pass per document, making it $O(N)$ for full-corpus
search---impractical beyond $\sim$1000 candidates.

\textbf{In production:} use both in a two-stage pipeline.  Stage 1: bi-encoder
(or hybrid BM25 + bi-encoder) retrieves top-100 candidates cheaply.  Stage 2:
cross-encoder re-ranks these 100 candidates to select the final top-5.  This
gives you the bi-encoder's efficiency at scale with the cross-encoder's
precision for the final selection.  The latency overhead of re-ranking 100
documents is $\sim$200--500ms, acceptable for most applications.

\textbf{Distillation:} for even better bi-encoder performance, use the
cross-encoder to score a large training set, then train the bi-encoder to
approximate these scores.

\redflags
\begin{itemize}
    \item Cannot explain why cross-encoders are higher quality (no mention
          of cross-attention or joint encoding)
    \item Proposes using a cross-encoder for full-corpus search
    \item Does not mention the two-stage pipeline as the standard pattern
    \item Confuses bi-encoder with cross-encoder architecturally
\end{itemize}

\interviewtest{Whether the candidate understands the fundamental
efficiency/quality tradeoff in information retrieval and can design a practical
two-stage system.  This is a core building block for any RAG system.}

\goodgreat
\begin{itemize}
    \item \badgeLfive{L5} Clearly explains both architectures, their tradeoffs, and
          the two-stage pipeline
    \item \badgeLsix{L6} Mentions ColBERT as a middle ground (late interaction),
          discusses distillation from cross-encoder to bi-encoder, and quantifies
          the quality gap
    \item \badgeLseven{L7} Discusses multi-vector representations, learned sparse
          models (SPLADE) as another point in the tradeoff space, and proposes
          adaptive re-ranking depth (more re-ranking for hard queries, less for
          easy ones)
\end{itemize}

\followups{What is ColBERT and where does it fall in the bi-encoder vs.\
cross-encoder spectrum?  How would you decide how many candidates to re-rank?
How does distillation from cross-encoder to bi-encoder work?}
\end{interviewq}

\bigskip

%% QUESTION 9 %%
\begin{interviewq}
\textbf{Q: Your RAG system is returning correct documents but the LLM's
answers are still wrong. What would you investigate?}

\badgeLsix{L6} \quad \badgetype{Debugging} \quad \badgefreq{\freqmed}

\stronganswer
If retrieval is correct but answers are wrong, the problem is in the
\textbf{generation layer}.  I would investigate in order:

(1)~\textbf{Prompt engineering:} Is the system prompt clearly instructing the
model to use the retrieved context?  Are there explicit grounding instructions
(``Answer based only on the provided documents'')?  Are chunks formatted
clearly with source labels?

(2)~\textbf{Context window overflow:} Are we stuffing too many chunks,
causing the model to miss the relevant one?  The ``lost in the middle''
effect may be burying the answer.  Try fewer chunks or reorder them.

(3)~\textbf{Chunk quality:} Even though the right \emph{document} is
retrieved, does the specific \emph{chunk} contain sufficient information?
The answer may span a chunk boundary.  Try larger chunks or the parent-child
pattern.

(4)~\textbf{Model limitations:} The model may lack the reasoning capability to
synthesize information from multiple chunks.  Try a more capable model or
decompose the query into sub-queries.

(5)~\textbf{Conflicting information:} Multiple retrieved chunks may contain
contradictory information (e.g., outdated vs.\ current docs).  Add metadata
filtering (prefer recent documents) or conflict resolution instructions.

(6)~\textbf{Instruction drift:} In multi-turn conversations, the grounding
instruction may be pushed out of the model's effective attention window.
Re-inject grounding instructions periodically.

\redflags
\begin{itemize}
    \item Immediately blames the retrieval system without acknowledging the
          problem statement
    \item Cannot enumerate specific generation-layer failure modes
    \item Proposes only ``use a better model'' without actionable diagnostics
\end{itemize}

\interviewtest{Whether the candidate can systematically debug a complex system
with interacting components.  This tests production engineering mindset---
knowing how to decompose and diagnose, not just build.}

\goodgreat
\begin{itemize}
    \item \badgeLfive{L5} Identifies prompt engineering and context ordering as likely
          issues, proposes concrete fixes
    \item \badgeLsix{L6} Systematically works through the generation failure taxonomy,
          discusses the lost-in-the-middle effect, and proposes A/B testing changes
    \item \badgeLseven{L7} Designs a diagnostic pipeline: compare model output with
          and without context (to detect if context is being ignored), instrument
          attention patterns, and propose fine-tuning for grounding faithfulness
\end{itemize}

\followups{How would you measure whether the model is actually using the
retrieved context vs.\ its parametric knowledge?  How would you instrument
a RAG system for debugging?  What metrics would you track to detect this
failure mode automatically?}
\end{interviewq}

\bigskip

%% QUESTION 10 %%
\begin{interviewq}
\textbf{Q: How would you scale a RAG system to handle 10 million documents and
1000 queries per second?}

\badgeLsix{L6} \quad \badgetype{System Design} \quad \badgefreq{\freqmed}

\stronganswer
At this scale, every component needs careful engineering.

\textbf{Indexing:} 10M documents at 512 tokens/chunk, assume 5 chunks per
document $=$ 50M chunks.  At 1024-dimensional embeddings in float32, that is
$50\text{M} \times 1024 \times 4$ bytes $\approx 200$\,GB of vectors.  Use
HNSW with scalar quantization (int8) to reduce to $\sim$50\,GB, fitting in
memory on a single high-memory machine, or shard across multiple nodes.

\textbf{Retrieval at 1K QPS:} HNSW queries take $\sim$1--5ms each.
A single machine can handle 200--1000 QPS depending on index size and parameters.
For 1K QPS, deploy 2--4 retrieval replicas behind a load balancer.  BM25 via
Elasticsearch scales horizontally with sharding.

\textbf{Re-ranking:} Cross-encoder re-ranking of 50 candidates takes
$\sim$200ms on GPU.  At 1K QPS, need $\sim$200 GPU-seconds per second;
deploy a pool of $\sim$4$-$8 GPU instances with batched inference.

\textbf{LLM generation:} The bottleneck.  At 1K QPS with $\sim$500ms per
generation, need $\sim$500 concurrent generation slots.  Use a high-throughput
serving framework (vLLM, TensorRT-LLM) with continuous batching.

\textbf{Caching:} Cache embeddings for frequent queries.  Cache full responses
for repeated identical queries (with TTL for freshness).  Semantic caching for
near-duplicate queries.

\textbf{Optimization:} Use Matryoshka embeddings to search at lower dimensions
for initial candidates, then re-score at full dimension.  Asynchronous indexing
pipeline for document updates.

\redflags
\begin{itemize}
    \item Cannot estimate memory requirements for the vector index
    \item Does not mention horizontal scaling or replication
    \item Ignores the LLM generation bottleneck
    \item No mention of caching
\end{itemize}

\interviewtest{Whether the candidate can reason about production infrastructure
at scale, performing back-of-the-envelope calculations and identifying
bottlenecks.  This separates ML engineers from ML scientists.}

\goodgreat
\begin{itemize}
    \item \badgeLfive{L5} Identifies the key bottlenecks and proposes horizontal
          scaling for retrieval
    \item \badgeLsix{L6} Provides concrete back-of-envelope calculations for memory,
          QPS per machine, and GPU requirements; mentions caching and batching
    \item \badgeLseven{L7} Discusses index sharding strategies, consistency
          guarantees during index updates, multi-region deployment for latency,
          and cost optimization (quantized embeddings, tiered retrieval with
          cheaper pre-filtering)
\end{itemize}

\followups{How would you handle real-time document updates while serving
queries?  What consistency guarantees do you need?  How would you shard the
vector index across multiple machines?}
\end{interviewq}


% ============================================================================
% CHAPTER SUMMARY
% ============================================================================
\section*{Chapter Summary}

\begin{enumerate}
    \item \textbf{RAG architecture} combines retrieval and generation to ground
          LLM responses in external knowledge, solving the knowledge cutoff,
          hallucination, and private data access problems.

    \item \textbf{Retrieval paradigms} span a quality-efficiency spectrum:
          BM25 (sparse, fast, lexical) $\to$ bi-encoder (dense, semantic) $\to$
          ColBERT (late interaction, high quality) $\to$ cross-encoder (highest
          quality, re-ranking only).  Production systems use hybrid retrieval
          with RRF fusion.

    \item \textbf{Chunking} is deceptively important: chunk size, overlap, and
          method should be treated as hyperparameters.  The parent-child
          pattern resolves the precision-context tradeoff.

    \item \textbf{Vector databases} use ANN algorithms (HNSW, IVF, PQ) to
          make dense retrieval practical at scale, with tradeoffs between
          recall, latency, and memory.

    \item \textbf{Re-ranking} with cross-encoders or LLMs dramatically improves
          precision.  The standard pattern is retrieve-many-then-re-rank-to-few.

    \item \textbf{Advanced patterns}---HyDE, query decomposition, multi-hop
          retrieval, Self-RAG, agentic RAG, and context compression---address
          the limitations of basic retrieve-then-generate.

    \item \textbf{RAG evaluation} is a three-layer problem: retrieval quality
          (Recall@$K$, MRR, NDCG), generation quality (relevance, completeness),
          and faithfulness (does the answer use only retrieved evidence?).

    \item \textbf{RAG vs.\ fine-tuning vs.\ long context} are complementary
          tools: RAG for dynamic knowledge with attribution, fine-tuning for
          behavior change, long context for small-corpus analysis.  The most
          common production pattern combines fine-tuning with RAG.
\end{enumerate}

\bigskip
\begin{keyinsight}[Chapter Summary---What Interviewers Want]
RAG is the most commonly tested system design topic for applied ML roles
building LLM-powered products.  The three things that distinguish a strong
candidate are: (1)~\textbf{Full pipeline thinking}---not just ``use an
embedding model and an LLM,'' but a complete design covering chunking,
hybrid retrieval, re-ranking, prompt construction, evaluation, and monitoring.
(2)~\textbf{Retrieval depth}---understanding the bi-encoder vs.\ cross-encoder
tradeoff, hybrid retrieval, and why BM25 remains relevant.
(3)~\textbf{Diagnostic ability}---when the system fails, knowing how to
decompose the failure into retrieval, ranking, and generation layers and
systematically fix each.  Candidates who treat RAG as a black box will not
pass; candidates who can design, evaluate, debug, and scale a RAG system
will earn strong hires.
\end{keyinsight}
